\documentclass[14pt,a4paper,oneside]{extbook}
\usepackage[utf8]{vietnam}
\usepackage{amsmath,amssymb,mathrsfs,fancyhdr,enumerate,multirow,makecell}
\usepackage{tkz-euclide,tikz-3dplot,tikz,tkz-tab}
\usetikzlibrary{arrows.meta,calc,intersections,angles,quotes,backgrounds,shapes.geometric,patterns,through,shadows,positioning}
\usepackage{yhmath}
\usepackage[top=0.7cm, bottom=0.7cm, left=1.0cm, right=1.0cm, headheight=16pt, headsep=3mm, footskip=5mm]{geometry}
\usepackage[hidelinks,unicode]{hyperref}
\usepackage{qrcode}
\usepackage[many]{tcolorbox}
\usepackage{fontawesome}

\pagestyle{fancy}
\fancyhf{}
\fancyhead[L]{\footnotesize\bfseries\color{teal!85!black}\faBook~SỔ TAY KIẾN THỨC TOÁN 5}
\fancyhead[R]{\footnotesize\bfseries\color{blue!85!black}Lớp Toán Thầy Lĩnh -- loptoanthaylinh.vn}
\fancyfoot[C]{\footnotesize\bfseries\color{gray!80!black}-- Sổ tay kiến thức: Bài 39. Tìm hai số khi biết hiệu và tỉ số của hai số đó --}
\renewcommand{\headrulewidth}{0.4pt}

\begin{document}

% --- BANNER TIÊU ĐỀ BÀI HỌC ---
\noindent
\begin{tcolorbox}[
  enhanced,
  colback=teal!8!white,
  colframe=teal!80!black,
  arc=3mm,
  boxrule=1.1pt,
  width=\linewidth,
  halign=center,
  top=2.2mm, bottom=2.2mm,
  left=3mm, right=3mm,
  drop fuzzy shadow
]
  {\large\bfseries\color{teal!95!black}\faBookmark\ BÀI 39. TÌM HAI SỐ KHI BIẾT HIỆU VÀ TỈ SỐ}\\[2pt]
  {\scriptsize\color{black!75}\itshape Phương pháp giải toán Hiệu -- Tỉ: Hiệu số phần bằng nhau và bài toán tuổi tác kinh điển trên 1 trang A4}
\end{tcolorbox}

\vspace{3.2mm}
% --- KHỐI 1: BỐN BƯỚC GIẢI TOÁN HIỆU TỈ ---
\noindent
\begin{tcolorbox}[
  enhanced,
  title={\footnotesize\bfseries\color{white}\faKey\ 1. Bốn bước vàng giải bài toán Hiệu -- Tỉ},
  colback=blue!2!white,
  colframe=blue!75!black,
  arc=2.5mm, boxrule=0.8pt, top=1.8mm, bottom=1.8mm, left=2.5mm, right=2.5mm
]
  \footnotesize
  \textbf{Bước 1: Vẽ sơ đồ đoạn thẳng}\\
  $\bullet$ Số bé vẽ $a$ phần, Số lớn vẽ $b$ phần ($b > a$). Thể hiện đoạn chênh lệch chính là \textbf{Hiệu}.\\
  \textbf{Bước 2: Tìm hiệu số phần bằng nhau}\\
  \begin{tcolorbox}[colback=white, colframe=blue!75!black, boxrule=0.5pt, arc=1.5mm, halign=center, top=1mm, bottom=1mm]
    {\scriptsize\bfseries\color{blue!90!black}$\text{Hiệu số phần bằng nhau} = b - a\text{ (phần)}$}
  \end{tcolorbox}
  \textbf{Bước 3: Tìm giá trị của một phần:}\\
  $\bullet$ $\text{Giá trị } 1\text{ phần} = \text{Hiệu hai số} : \text{Hiệu số phần}$.\\
  \textbf{Bước 4: Tìm số bé và số lớn:}\\
  $\bullet$ $\text{Số bé} = \text{Giá trị } 1\text{ phần} \times a$ \quad | \quad $\text{Số lớn} = \text{Số bé} + \text{Hiệu}$ (hoặc Giá trị $1$ phần $\times b$).
\end{tcolorbox}

\vspace{3mm}

% --- KHỐI 2: MINH HỌA SƠ ĐỒ ĐOẠN THẲNG ---
\noindent
\begin{minipage}[t]{0.52\textwidth}
  \begin{tcolorbox}[
    enhanced,
    title={\footnotesize\bfseries\color{white}\faPencil\ Sơ đồ mẫu (Tỉ số 3/5, Hiệu 18)},
    colback=teal!2!white,
    colframe=teal!75!black,
    arc=2.5mm, boxrule=0.8pt, top=1.8mm, bottom=1.8mm, left=2mm, right=2mm,
    halign=center
  ]
    \begin{tikzpicture}[scale=0.68]
      \draw[blue!85!black, thick] (0, 1) node[left, font=\tiny] {Số bé:} -- (3, 1);
      \draw (0, 0.9)--(0, 1.1); \draw (1, 0.9)--(1, 1.1); \draw (2, 0.9)--(2, 1.1); \draw (3, 0.9)--(3, 1.1);
      \draw[red!85!black, thick] (0, 0) node[left, font=\tiny] {Số lớn:} -- (5, 0);
      \draw (0, -0.1)--(0, 0.1); \draw (1, -0.1)--(1, 0.1); \draw (2, -0.1)--(2, 0.1); \draw (3, -0.1)--(3, 0.1); \draw (4, -0.1)--(4, 0.1); \draw (5, -0.1)--(5, 0.1);
      \draw[dashed] (3, 1) -- (3, 0);
      \draw[decorate, decoration={brace, amplitude=3pt}, thick] (5, -0.2) -- (3, -0.2) node[midway, below=3pt, font=\scriptsize\bfseries] {18};
    \end{tikzpicture}\\[2pt]
    {\tiny\color{teal!95!black}\bfseries Số bé $= 18 : 2 \times 3 = 27$; Số lớn $= 27 + 18 = 45$}
  \end{tcolorbox}
\end{minipage}%
\hfill
\begin{minipage}[t]{0.46\textwidth}
  \begin{tcolorbox}[
    enhanced,
    title={\footnotesize\bfseries\color{white}\faBolt\ Công thức gộp tính nhanh},
    colback=purple!2!white,
    colframe=purple!70!black,
    arc=2.5mm, boxrule=0.8pt, top=1.8mm, bottom=1.8mm, left=2mm, right=2mm
  ]
    \footnotesize
    \begin{tcolorbox}[colback=white, colframe=purple!70!black, boxrule=0.5pt, arc=1.5mm, halign=center, top=1mm, bottom=1mm]
      {\scriptsize\bfseries\color{purple!90!black}$\text{Số bé} = \dfrac{\text{Hiệu}}{b - a} \times a$}\\[1.5pt]
      {\scriptsize\bfseries\color{blue!85!black}$\text{Số lớn} = \text{Số bé} + \text{Hiệu}$}
    \end{tcolorbox}
    Hai số tìm được có hiệu chính xác là đề bài.
  \end{tcolorbox}
\end{minipage}

\vspace{3mm}

% --- KHỐI 3: SƠ ĐỒ TƯ DUY ---
\noindent
\begin{tcolorbox}[
  enhanced,
  title={\footnotesize\bfseries\color{white}\faSitemap\ Sơ đồ tư duy: Giải bài toán Hiệu -- Tỉ},
  colback=yellow!2!white,
  colframe=teal!85!black,
  arc=3mm, boxrule=0.9pt, top=1.8mm, bottom=1.8mm, left=1.5mm, right=1.5mm,
  halign=center
]
  \begin{tikzpicture}[
    scale=0.92,
    every node/.style={transform shape},
    b1/.style={rectangle, rounded corners=4pt, fill=blue!15, draw=blue!70!black, text=blue!90!black, font=\bfseries\scriptsize, inner sep=4pt, align=center},
    b2/.style={rectangle, rounded corners=4pt, fill=teal!15, draw=teal!70!black, text=teal!90!black, font=\bfseries\scriptsize, inner sep=4pt, align=center},
    b3/.style={rectangle, rounded corners=4pt, fill=orange!15, draw=orange!70!black, text=orange!90!black, font=\bfseries\scriptsize, inner sep=4pt, align=center},
    b4/.style={rectangle, rounded corners=4pt, fill=green!15, draw=green!70!black, text=green!90!black, font=\bfseries\scriptsize, inner sep=4pt, align=center}
  ]
    \node[b1] (s1) at (-6.3, 0) {1. Vẽ sơ đồ\\đoạn thẳng};
    \node[b2] (s2) at (-2.1, 0) {2. Hiệu số phần\\($b - a$)};
    \node[b3] (s3) at (2.1, 0) {3. Giá trị\\$1$ phần};
    \node[b4] (s4) at (6.3, 0) {4. Tìm số bé\\và số lớn};

    \draw[teal!85!black, thick, -{Stealth[length=2.5mm]}] (s1) -- (s2);
    \draw[teal!85!black, thick, -{Stealth[length=2.5mm]}] (s2) -- (s3);
    \draw[teal!85!black, thick, -{Stealth[length=2.5mm]}] (s3) -- (s4);
  \end{tikzpicture}
\end{tcolorbox}

\vspace{3mm}

% --- KHỐI 4: TOÁN TUỔI TÁC BẤT BIẾN & BẪY TRẮC NGHIỆM ---
\noindent
\begin{minipage}[t]{0.54\textwidth}
  \begin{tcolorbox}[
    enhanced,
    title={\footnotesize\bfseries\color{white}\faUsers\ Bài toán tuổi tác: Hiệu số tuổi KHÔNG ĐỔI},
    colback=green!3!white,
    colframe=green!60!black,
    arc=2.5mm, boxrule=0.8pt, top=1.8mm, bottom=1.8mm, left=2.2mm, right=2.2mm
  ]
    \footnotesize
    \textbf{Nguyên lý vàng}: Qua mỗi năm, mỗi người đều thêm $1$ tuổi, do đó \textbf{Hiệu số tuổi giữa hai người luôn không thay đổi theo thời gian}!\\
    $\triangleright$ Bố hơn con $28$ tuổi thì $5$ năm trước hay $10$ năm sau bố vẫn hơn con $28$ tuổi.
  \end{tcolorbox}
\end{minipage}%
\hfill
\begin{minipage}[t]{0.44\textwidth}
  \begin{tcolorbox}[
    enhanced,
    title={\footnotesize\bfseries\color{white}\faExclamationTriangle\ Mẹo tránh bẫy trắc nghiệm},
    colback=red!2!white,
    colframe=red!70!black,
    arc=2.5mm, boxrule=0.8pt, top=1.8mm, bottom=1.8mm, left=2.2mm, right=2.2mm
  ]
    \footnotesize
    $\bullet$ \textbf{Lỗi nhầm sang Tổng -- Tỉ}: Ở bài toán Hiệu -- Tỉ, bước 2 phải làm phép \textbf{trừ} ($b - a$), rất nhiều bạn quen tay làm phép cộng!\\
    $\bullet$ Tìm số lớn: Lấy số bé \textbf{cộng} với Hiệu.
  \end{tcolorbox}
\end{minipage}

\end{document}
